\documentclass[11pt]{article}
\usepackage[margin=1.1in]{geometry}
\usepackage{amsmath}
\usepackage{booktabs}
\usepackage{xcolor}
\usepackage[colorlinks=true,linkcolor=blue!60!black,citecolor=green!40!black,urlcolor=blue!60!black]{hyperref}

\title{Reviewing Documents with a Coding Agent}
\author{The pdfloop example paper}
\date{}

\begin{document}
\maketitle

\begin{abstract}
Writing a paper with a coding agent is fast until the review step. The author reads the
compiled PDF, but the agent edits the source, so every comment has to be copied out of the
PDF, located in the source and explained again in a chat window. This example paper exists
to be reviewed with pdfloop: open it, select a sentence, leave a comment, and send it to
your agent.
\end{abstract}

\section{Introduction}
\label{sec:intro}

Most technical documents are written in a markup language and read as a PDF. LaTeX
\cite{lamport1994latex} and Typst \cite{typst} both separate the two: the author edits
plain text, a compiler produces pages. Coding agents are good at the first half. They can
restructure a section, fix notation everywhere, or rewrite a paragraph in a different tone.
They cannot see what the reader sees.

The review loop closes that gap. The reader comments on the rendered page, the agent
receives each comment with the quoted text and its surroundings, edits the source,
rebuilds, and replies. Section~\ref{sec:loop} describes the loop and
Section~\ref{sec:cost} estimates what it saves.

\section{The review loop}
\label{sec:loop}

A review round has four steps. The reader selects text and writes a comment. The reader
sends all open comments at once. The agent applies them and rebuilds the PDF. The viewer
reloads the new PDF at the same position, and each comment shows the agent's reply.

The viewer never touches the project. It only records comments and requests; the agent
decides what to change, and the reader can ask it to commit when the round is done.
Table~\ref{tab:roles} lists who does what.

\begin{table}[h]
\centering
\begin{tabular}{lll}
\toprule
Step & Reader & Agent \\
\midrule
Comment & selects text, writes a note & \\
Send & presses Cmd+Enter & wakes up \\
Apply & & edits source, rebuilds \\
Reply & reads the reply in place & writes one reply per comment \\
\bottomrule
\end{tabular}
\caption{Division of work in one review round.}
\label{tab:roles}
\end{table}

\section{What it saves}
\label{sec:cost}

Let $n$ be the number of comments in a round, $t_c$ the time to write one comment, and
$t_l$ the time to locate the commented text in the source and describe it to the agent.
Reviewing in a chat window costs
\begin{equation}
T_{\text{chat}} = n\,(t_c + t_l),
\label{eq:chat}
\end{equation}
while reviewing on the page costs $T_{\text{page}} = n\,t_c$, because the quote and its
context travel with the comment. For a round of twenty comments with $t_l = 30$\,s,
Equation~\eqref{eq:chat} gives ten minutes spent only on pointing at text.

\section{Conclusion}

Comment where you read, edit where you write, and let the agent move between the two.
PDF.js \cite{pdfjs} renders the pages; the rest is a small Python server and a
skill file that tells the agent how to take part.

\begin{thebibliography}{9}
\bibitem{lamport1994latex} Leslie Lamport. \emph{\LaTeX: A Document Preparation System}.
Addison-Wesley, 2nd edition, 1994.
\bibitem{typst} Typst. \url{https://typst.app}.
\bibitem{pdfjs} Mozilla. PDF.js. \url{https://mozilla.github.io/pdf.js/}.
\end{thebibliography}

\end{document}
